\documentclass[11pt]{article}
\usepackage[margin=1in, headheight=14pt]{geometry}
\usepackage{amsmath, amssymb}
\usepackage{hyperref}
\usepackage{enumitem}
\usepackage{algorithm}
\usepackage{algpseudocode}
\usepackage{fancyhdr}
\usepackage{tcolorbox}
\usepackage{microtype}
\usepackage{tikz}
\usepackage{pgfplots}
\pgfplotsset{compat=1.18}
\usetikzlibrary{arrows.meta, positioning, calc, shapes.geometric}

\input{notation.tex}
\input{zk-notation.tex}
\input{environments.tex}
\input{lecture-macros.tex}

\pagestyle{fancy}
\fancyhf{}
\lhead{EECS 498/CSE 598: Zero-Knowledge Proofs}
\rhead{Winter 2026}
\lfoot{Lecture 22: Ecash, Blockchains, and Zerocash}
\rfoot{\thepage}
\renewcommand{\headrulewidth}{0.4pt}
\renewcommand{\footrulewidth}{0.4pt}

% Lecture-specific macros
\tcbuselibrary{skins, breakable}

\newtcolorbox{highlight}[1][]{
  colback=blue!5!white,
  colframe=blue!60!black,
  fonttitle=\bfseries,
  #1
}

\newtcolorbox{graybox}[1][]{
  colback=gray!8,
  colframe=gray!50,
  fonttitle=\bfseries,
  #1
}

\newcommand{\CRH}{\algo{CRH}}
\newcommand{\PRF}{\algo{PRF}}
\newcommand{\Com}{\algo{Com}}
\newcommand{\Zerocash}{\ensuremath{\mathsf{Zerocash}}}
\newcommand{\MT}{\mathsf{MT}}
\newcommand{\txmint}{\mathsf{tx}_{\mathrm{mint}}}
\newcommand{\txpour}{\mathsf{tx}_{\mathrm{pour}}}

\begin{document}

%-----------------------------------------------------------------------
% Title Block
%-----------------------------------------------------------------------
\begin{center}
  {\LARGE \textbf{EECS 498 / CSE 598: Zero-Knowledge Proofs}}\\[6pt]
  {\large Winter 2026}\\[4pt]
  {\large \textbf{Lecture 22: Ecash, Blockchains, and Zerocash}}\\[4pt]
  {\large Instructor: Paul Grubbs \quad
    \href{mailto:paulgrub@umich.edu}{\texttt{paulgrub@umich.edu}} \quad
    Beyster 4709}
\end{center}

\vspace{0.5em}
\hrule
\vspace{1em}

%-----------------------------------------------------------------------
% Agenda
%-----------------------------------------------------------------------
\section*{Agenda}
\begin{itemize}
  \item Chaumian ecash and blockchain background
  \item Sander--Ta Shma anonymous ecash
  \item Zerocash coins and pour transactions
  \item Other applications: zkVMs, rollups, and private smart contracts
\end{itemize}

\hrule
\vspace{1em}

%-----------------------------------------------------------------------
\section{Chaumian Ecash and Blockchain Background}
%-----------------------------------------------------------------------

Electronic cash systems combine two basic goals:
\begin{enumerate}
  \item value conservation, so users cannot create money by deviating from the
        protocol;
  \item payment privacy, so a merchant or bank does not automatically learn
        the spender's identity from an ordinary payment transcript.
\end{enumerate}

Chaum's blind-signature construction is the canonical starting point.  A user
chooses a random serial number $sn$, blinds it, and sends the blinded value to
the bank.  After deducting the coin's value from the user's account, the bank
signs the blinded serial number.  The user unblinds the response and obtains a
signature $\sigma_{sn}$ on the original serial number.

Spending is simple:
\begin{enumerate}
  \item the user sends $(sn,\sigma_{sn})$ to the merchant;
  \item the merchant asks the bank whether $sn$ has already appeared;
  \item the bank rejects repeated serial numbers and honors fresh ones.
\end{enumerate}

\begin{remark}
Blind signatures hide the serial number during withdrawal, so the bank cannot
link the later spend transcript to the withdrawal transcript.  At the same
time, the bank remains a trusted online party and its signing key is a single
point of failure.
\end{remark}

\begin{definition}[Blockchain ledger]
A blockchain is an append-only distributed ledger that gives all participants a
common view of previously accepted transactions.  The ledger state acts as a
public record of ownership changes, spent objects, and auxiliary consensus
data.
\end{definition}

The blockchain viewpoint replaces the online bank by a replicated public state.
For payment systems built on top of a ledger, the main validity checks become:
\begin{itemize}
  \item whether the spent object previously appeared on the ledger,
  \item whether its serial number is fresh,
  \item whether the transaction preserves value.
\end{itemize}

This shift turns anonymous payments into a proof problem about public state:
membership in a ledger accumulator, ownership of a coin witness, and validity
of a value-conservation relation.

%-----------------------------------------------------------------------
\section{Sander--Ta Shma Anonymous Ecash}
%-----------------------------------------------------------------------

Sander--Ta Shma removes the bank's signing key and replaces signed coins by
commitments recorded in a public Merkle tree.  The bank's role is reduced to
maintaining the integrity of that public database.

\subsection{Coin Creation and Spending}

A Pedersen-style commitment is used.  A user with identity
$id$ samples
\[
  sn, r \getsr \F
\]
and computes
\[
  C = sn \cdot G + rH.
\]
The bank inserts $C$ into the public Merkle tree of issued coins.  The spend
relation then consists of revealing the serial number $sn$ together with:
\begin{itemize}
  \item a proof that $C$ lies on a Merkle path to the current root $\mathsf{rt}$;
  \item a proof of knowledge that $C$ opens to $(sn,r)$.
\end{itemize}

\begin{proposition}
If the commitment opening proof is sound and the Merkle-path proof is sound,
then an accepted spend corresponds to some committed coin in the bank's public
database.
\end{proposition}

\begin{proof}[Proof sketch]
The path proof shows that some commitment $C$ is a leaf of the Merkle tree with
root $\mathsf{rt}$.  The opening proof shows that the same $C$ is consistent
with the revealed serial number $sn$.  Therefore an accepted payment is linked
to some previously issued committed coin, even though the zero-knowledge proof
hides which leaf belongs to the spender.
\end{proof}

\subsection{Double-Spend Identification}

The extended construction adds traceability for double spends.  During coin
creation, the user chooses values $v_1,v_2,sn,r$ and forms
\[
  C = v_1 G_1 + v_2 G_2 + sn G_3 + rH,
\]
subject to the relation
\[
  v_2 = v_1 - id.
\]
Equivalently,
\[
  id = v_1 - v_2.
\]
The user proves knowledge of $(v_1,v_2,sn,r)$ and proves that this linear
relation holds.

To spend the coin at a merchant, the merchant derives a challenge
\[
  \mathit{chal} = H(\text{merchant ID}, \text{timestamp})
\]
and the user reveals
\[
  V = v_1 + \mathit{chal} \cdot v_2
\]
together with the serial number $sn$ and the usual path/opening proofs.

\begin{graybox}[title={Double-Spend Tracing Rule}]
If the same serial number is spent twice under distinct challenges
$\mathit{chal}_1 \neq \mathit{chal}_2$, then
\[
  V_1 = v_1 + \mathit{chal}_1 v_2
  \qquad\text{and}\qquad
  V_2 = v_1 + \mathit{chal}_2 v_2.
\]
These two equations determine
\[
  v_2 = \frac{V_1 - V_2}{\mathit{chal}_1 - \mathit{chal}_2},
  \qquad
  v_1 = V_1 - \mathit{chal}_1 v_2,
  \qquad
  id = v_1 - v_2.
\]
\end{graybox}

\begin{proposition}
Two accepted spends of the same coin under distinct merchant challenges reveal
the spender's identity.
\end{proposition}

\begin{proof}
Subtracting the two spend equations yields
\[
  V_1 - V_2 = (\mathit{chal}_1 - \mathit{chal}_2) v_2.
\]
Because the challenges are distinct, $\mathit{chal}_1 - \mathit{chal}_2$ is
invertible, so $v_2$ is determined uniquely.  Substituting back gives $v_1$,
and then $id = v_1 - v_2$.
\end{proof}

\begin{remark}
Ordinary one-time spending remains anonymous because one transcript reveals
only the single linear form $V = v_1 + \mathit{chal} v_2$.  Traceability
appears only after two distinct challenges are answered for the same serial
number.
\end{remark}

%-----------------------------------------------------------------------
\section{Zerocash}
%-----------------------------------------------------------------------

\Zerocash{} moves anonymous payments to a blockchain setting.  The public
ledger stores coin commitments and spent serial numbers, while zero-knowledge
proofs justify minting and spending without exposing the underlying private
coin data.

\subsection{Addresses, Coins, and Ledger State}

Each address has a secret key $a_{sk}$ and a public key
\[
  a_{pk} = \PRF_{a_{sk}}(0).
\]
A coin is described by the tuple
\[
  c = (a_{pk}, v, \rho, r, s, cm),
\]
where:
\begin{itemize}
  \item $v$ is the coin value;
  \item $\rho$ is serial-number randomness;
  \item $r$ is commitment randomness for the address layer;
  \item $s$ is commitment randomness for the value layer;
  \item $cm$ is the final coin commitment.
\end{itemize}

The note commitment is built in two layers:
\[
  k = \Com_r(a_{pk}\|\rho),
  \qquad
  cm = \Com_s(v\|k).
\]
The ledger maintains:
\begin{itemize}
  \item a Merkle tree over all coin commitments $cm$;
  \item a public list of spent serial numbers.
\end{itemize}

\subsection{Minting}

To mint a coin of value $v$ for address $a_{pk}$, sample
\[
  \rho, r, s \getsr \bits^\lambda
\]
and compute
\[
  k = \Com_r(a_{pk}\|\rho),
  \qquad
  cm = \Com_s(v\|k).
\]
The mint transaction publishes the commitment data needed for ledger validity,
while the ownership secrets remain hidden.  A convenient explicit form is
\[
  \txmint = (cm, v, k, s),
\]
which makes the algebraic structure of the note commitment explicit.

\begin{remark}
The commitment tree plays the same structural role as the tree in
Sander--Ta Shma: it compresses the set of minted coins into a short public root
that can be used inside zero-knowledge membership proofs.
\end{remark}

\subsection{Pour Transactions}

A simplified pour consumes one old coin and creates two new
coins.  If the consumed coin is owned under secret key $a_{sk}^{old}$ and has
serial randomness $\rho^{old}$, its serial number is computed as
\[
  sn = \PRF_{a_{sk}^{old}}(\rho^{old}).
\]
Two new recipient addresses
\[
  a_{pk,1}^{new},
  \qquad
  a_{pk,2}^{new}
\]
receive two freshly sampled notes with commitments
\[
  cm_1^{new}, cm_2^{new}.
\]
The central invariant is value conservation:
\[
  v^{old} = v_1^{new} + v_2^{new}.
\]

The internal steps are:
\begin{enumerate}
  \item compute the old serial number $sn$ from $(a_{sk}^{old}, \rho^{old})$;
  \item sample fresh randomness $\rho_i^{new}, r_i^{new}, s_i^{new}$ for
        $i \in \{1,2\}$;
  \item compute the new layered commitments
        \[
          k_i^{new} = \Com_{r_i^{new}}(a_{pk,i}^{new}\|\rho_i^{new}),
          \qquad
          cm_i^{new} = \Com_{s_i^{new}}(v_i^{new}\|k_i^{new});
        \]
  \item generate a zero-knowledge proof that the old note is in the ledger
        tree, is owned by the spender, and has been re-expressed as the two new
        notes without violating value conservation.
\end{enumerate}

\begin{highlight}[title={Public Inputs in a Simplified Pour Proof}]
The public statement checked by the ledger contains the data needed for
uniqueness and ledger consistency:
\begin{itemize}
  \item the Merkle-tree root $\mathsf{rt}$ used to authenticate the consumed note;
  \item the revealed serial number $sn$ of the consumed note;
  \item the new note commitments $cm_1^{new}$ and $cm_2^{new}$;
  \item any public balance or auxiliary transaction data carried by the ledger.
\end{itemize}
\end{highlight}

\begin{proposition}
Assume the pour proof is sound, the ledger rejects repeated serial numbers, and
the proof enforces membership and value conservation.  Then an accepted pour
transaction consumes an existing note exactly once and re-expresses its value
as valid new output notes.
\end{proposition}

\begin{proof}[Proof sketch]
Soundness of the proof ties the revealed serial number to a committed old note
that lies in the Merkle tree under the public root $\mathsf{rt}$.  Serial
number freshness prevents reusing the same consumed note in a later accepted
transaction.  The value-conservation constraint forces the sum of output values
to match the value carried by the consumed note.  Therefore acceptance implies
one-time consumption of existing value together with correct creation of the new
notes.
\end{proof}

\begin{remark}
The public ledger reveals the serial number of the spent note and the new note
commitments, but not the spender's address secret key, the old note's address,
or the private values and randomness used to form the new notes.
\end{remark}

%-----------------------------------------------------------------------
\section{Other Applications}
%-----------------------------------------------------------------------

The same proof patterns reappear in several later systems:
\begin{itemize}
  \item \emph{zkVMs} prove correct execution of a virtual machine while
        revealing only the claimed public input and output;
  \item \emph{rollups} aggregate many ledger transitions into a compact proof
        checked by a base chain;
  \item \emph{private smart contracts} combine hidden state with proofs that
        each state transition satisfies the contract rules.
\end{itemize}

All three settings rely on the same ingredients developed throughout the later
part of the course: succinct proofs for state transitions, authenticated public
state, and carefully chosen representations of the hidden witness.

\bigskip
\hrule
\bigskip

\begin{thebibliography}{9}

\bibitem{Chaum82}
  D.~Chaum,
  ``Blind signatures for untraceable payments,''
  \emph{CRYPTO}, 1982.

\bibitem{ST99}
  T.~Sander and A.~Ta-Shma,
  ``Auditable, anonymous electronic cash,''
  \emph{CRYPTO}, 1999.

\bibitem{Zerocash14}
  E.~Ben-Sasson, A.~Chiesa, C.~Garman, M.~Green, I.~Miers, E.~Tromer,
  and M.~Virza,
  ``Zerocash: decentralized anonymous payments from Bitcoin,''
  \emph{IEEE Symposium on Security and Privacy}, 2014.

\end{thebibliography}

\end{document}
